\section{Reduction semantics}

The reduction semantics of configurations, written
$\conf{\rho}{P} \red_{s} \conf{\rho'}{P'}$ and indexed by the choice $s$ that
selects the redex, is the least relation closed under the rules of
\cref{tab:sts}. It follows the STS of VACCS: the first six rules are the
redexes, the last three close them under parallel composition, restriction
and structural congruence.

\begin{table}[h]
\centering
\footnotesize
$\begin{array}{l@{\quad}rcl@{\quad}l}
\textsc{Com} &
  \conf{\rho}{\tagged{t}{c!e} \parallel (\tagged{t'}{c?x}.P + K)} &\red_{\{t,t'\}}&
  \conf{\rho}{\mathbf{0} \parallel P[v/x]} & e \Downarrow v \\[2pt]
\textsc{Tau} &
  \conf{\rho}{\tagged{t}{\tau}.P + K} &\red_{t}& \conf{\rho}{P} \\[2pt]
\textsc{Op} &
  \conf{\rho}{\tagged{t}{U(\tilde q)}.P + K} &\red_{t}&
  \conf{\upd{U}}{P} & U : \mathsf{Op}(|\tilde q|) \\[2pt]
\textsc{Init} &
  \conf{\rho}{\tagged{t}{\initp}.P + K} &\red_{t}&
  \conf{\inirho}{P} \\[2pt]
\textsc{Meas} &
  \conf{\rho}{\tagged{t}{\M(\tilde q \rhd x)}.P + K} &\red_{t}&
  \conf{\br}{P[m/x]} & \M : \mathsf{Meas}(|\tilde q|),\ 0 \le m < |\M| \\[2pt]
\textsc{Rec} &
  \conf{\rho}{\tagged{t}{\rec{X}{P}}} &\red_{t}&
  \conf{\rho}{P[\tagged{t}{\rec{X}{P}}/X]} \\[2pt]
\textsc{Par} &
  \conf{\rho}{P \parallel Q} &\red_{s}& \conf{\rho'}{P' \parallel Q}
  & \conf{\rho}{P} \red_{s} \conf{\rho'}{P'} \\[2pt]
\textsc{Res} &
  \conf{\rho}{P \setminus c} &\red_{s}& \conf{\rho'}{P' \setminus c}
  & \conf{\rho}{P} \red_{s} \conf{\rho'}{P'} \\[2pt]
\textsc{Cong} &
  \conf{\rho}{P} &\red_{s}& \conf{\rho'}{Q}
  & P \equiv P_1,\ \conf{\rho}{P_1} \red_{s} \conf{\rho'}{Q_1},\ Q_1 \equiv Q
\end{array}$
\caption{Reduction semantics of configurations.}
\label{tab:sts}
\end{table}

As in VACCS, the proof of the Harmony lemma goes through the \emph{shapes}
of the transitions: up to congruence, a process that performs an action is a
redex under some restrictions, in parallel with a rest. Here
$(\nu \tilde d)\,P$ stands for $P \setminus d_1 \cdots \setminus d_n$, with
$\tilde d$ fresh for everything but $P$ (by $\alpha$-conversion).

\begin{lemma}[Shapes of the transitions]\label{lem:shapes}
In the program LTS,
\begin{enumerate}
\item if $P \step{c!v}{s} P'$ then $s = t$ and
  $P \equiv (\nu \tilde d)(\tagged{t}{c!e} \parallel R)$,
  $P' \equiv (\nu \tilde d)(\mathbf{0} \parallel R)$ for some $t, e, R, \tilde d$
  with $e \Downarrow v$ and $c \notin \tilde d$;
\item if $P \step{c?v}{s} P'$ then $s = t$ and
  $P \equiv (\nu \tilde d)((\tagged{t}{c?x}.P_1 + K) \parallel R)$,
  $P' \equiv (\nu \tilde d)(P_1[v/x] \parallel R)$ for some
  $t, x, P_1, K, R, \tilde d$ with $c \notin \tilde d$;
\item if $\conf{\rho}{K} \step{\tau}{s} \conf{\rho'}{P'}$ for a sum $K$, then
  $K \equiv \alpha.P + K'$ and
  $\conf{\rho}{\alpha.P} \step{\tau}{s} \conf{\rho'}{P'}$ for some
  $\alpha, P, K'$.
\end{enumerate}
\end{lemma}

\begin{proof}
By induction on the derivation of the transition. A conditional is removed by
\textsc{IteT} or \textsc{IteF}; for ParL and ParR, the other component is
pushed into the rest $R$ under the restrictions by \textsc{ParComm},
\textsc{ParAssoc} and \textsc{ResPar} (after $\alpha$-conversion of $\tilde d$);
for Restr, the restriction is added to $\tilde d$; for SumL and SumR, the
summand that moves is brought in front by \textsc{SumComm} and
\textsc{SumAssoc}.
\end{proof}

\begin{lemma}[Harmony]\label{lem:harmony}
For all $\rho, P, s, \rho', Q$,
\[
\conf{\rho}{P} \red_{s} \conf{\rho'}{Q}
\iff
\conf{\rho}{P} \step{\tau}{s} \conf{\rho'}{Q'} \text{ for some } Q' \equiv Q.
\]
In particular, for a parallel composition,
$\conf{\rho}{P \parallel Q} \red_{s} \conf{\rho'}{R}$ iff
$\conf{\rho}{P \parallel Q} \step{\tau}{s} \conf{\rho'}{R'}$ for some
$R' \equiv R$.
\end{lemma}

\begin{proof}
($\Rightarrow$) By induction on the derivation of the reduction. The redexes
\textsc{Com}, \textsc{Tau}, \textsc{Op}, \textsc{Init}, \textsc{Meas} and
\textsc{Rec} are transitions of \cref{fig:conf}, the summand that moves being selected by
SumL. \textsc{Par} and \textsc{Res} are ParL and Restr. For \textsc{Cong}, let
$P \equiv P_1$ and, by induction hypothesis,
$\conf{\rho}{P_1} \step{\tau}{s} \conf{\rho'}{Q_1'}$ with $Q_1' \equiv Q_1$.
By \cref{lem:pre-harmony} this transition is a silent move of $P_1$, or a
request of $P_1$ synchronised with the memory; \cref{lem:cgr-bisim} gives the
same move of $P$, with the same choice and up to $\equiv$, and
\cref{lem:pre-harmony} again turns it into a transition of
$\conf{\rho}{P}$.

($\Leftarrow$) By \textsc{Cong} it is enough to show that
$\conf{\rho}{P} \step{\tau}{s} \conf{\rho'}{Q}$ implies
$\conf{\rho}{P} \red_{s} \conf{\rho'}{Q}$, by induction on the derivation of
the transition. For a sum, \cref{lem:shapes}(3) brings the summand that moves
in front, and the redex is \textsc{Tau}, \textsc{Op}, \textsc{Init} or
\textsc{Meas}. IteT
and IteF are \textsc{Cong} with \textsc{IteT} and \textsc{IteF}; ParL is
\textsc{Par}; ParR is \textsc{Par} under \textsc{ParComm}; Restr is
\textsc{Res}; Rec is \textsc{Rec}. For SynchL, with
$P \step{c!v}{t} P'$ and $Q \step{c?v}{t'} Q'$, \cref{lem:shapes}(1) and (2)
give $P \parallel Q \equiv
(\nu \tilde d)(\tagged{t}{c!e} \parallel R) \parallel
(\nu \tilde d')((\tagged{t'}{c?x}.P_1 + K) \parallel R')$; the restrictions
$\tilde d$ and $\tilde d'$ are extruded one by one to the top with
\textsc{ResPar}, the output and the input are brought side by side with
\textsc{ParComm} and \textsc{ParAssoc}, and the redex is \textsc{Com}, with
choice $\{t,t'\}$. SynchR is SynchL under \textsc{ParComm}.
\end{proof}

\iffalse
The proof of the Harmony lemma goes through a restatement at the level of
programs, which is the form used in the next section. Write
$\rho \xrightarrow{\lambda} \rho'$ for the effect of an internal action on the
memory: $\rho \xrightarrow{\tau} \rho$, $\rho \xrightarrow{\req{U}} \upd{U}$,
$\rho \xrightarrow{\initp} \inirho$ and $\rho \xrightarrow{\mrecv{m}} \br$.

\begin{lemma}[Harmony, program level]\label{lem:harmony-prog}
$\conf{\rho}{P} \red_{s} \conf{\rho'}{Q}$ iff $P \step{\lambda}{s} Q'$ and
$\rho \xrightarrow{\lambda} \rho'$ for some internal action $\lambda$ and some
$Q' \equiv Q$.
\end{lemma}
\fi
